% GENERATED FILE. Edit canonical lessons, then run npm run generate:books.
% canonical-source-hash: fnv1a64:dce1e31ebe7095e4
% canonical-lessons: TE-C449-addanki, TE-C449-mallimpu, TE-C449-nilicipovu, TE-C449-stambhincu, TE-C449-musukupovu

\chapter{Roads Blocked}
\label{ch:te-roads-blocked}
\hlchaptermodality{449}

\begin{chapteropening}
\textbf{By the end of this chapter:} \emph{I can say an obstacle, a diversion, to stop dead, to seize up and to get blocked.}

Complete the last lesson of chapter 449: I can say an obstacle, a diversion, to stop dead, to seize up and to get blocked.
\end{chapteropening}

\section[\texorpdfstring{\te{అడ్డంకి} (aḍḍaṅki)}{అడ్డంకి (aḍḍaṅki)}]{\te{అడ్డంకి} (aḍḍaṅki) — an obstacle; a hindrance}

\label{lesson:TE-C449-addanki}



\begin{quote}
Before the new one: say the Telugu for to take a turn, then the Telugu for a shortcut.
\end{quote}

\subsection*{You'll want to know: \te{అడ్డంకి}}
\textbf{\te{అడ్డంకి}} — \emph{aḍḍaṅki} — \textquotedblleft{}an obstacle; a hindrance\textquotedblright{}.

\begin{quote}
\te{రోడ్డు} \te{మీద} \te{పడిన} \te{చెట్టు} \te{పెద్ద} \te{అడ్డంకి} \te{అయింది}.
\end{quote}

\begin{quote}
\emph{rōḍḍu mīda paḍina ceṭṭu pedda aḍḍaṅki ayindi.}
\end{quote}

\begin{quote}
The tree that fell on the road became a big obstacle.
\end{quote}

\te{అడ్డు}, a hindrance, + -\te{ంకి}. It goes with plans too: \te{పనికి} \te{అడ్డంకి}, an obstacle to the work.

\begin{cousinweb}
\textbf{\te{అడ్డంకి}} (\emph{aḍḍaṅki}) is native Dravidian, built on \te{అడ్డు}, \textquotedblleft{}to hinder; a hindrance, a cross-bar\textquotedblright{} (DEDR 83), with Tamil \emph{aṭai}, \textquotedblleft{}to block\textquotedblright{}.
\end{cousinweb}

\begin{grammarlens}[title={What you've built}]
One more word for getting around.
\end{grammarlens}

\subsection*{Your turn}
\begin{itemize}
\raggedright
  \item \emph{Say it:} \emph{aḍḍaṅki}
  \item \emph{Say it:} \emph{aḍḍaṅki}, once more
  \item \emph{Say it:} \emph{rōḍḍu mīda paḍina ceṭṭu pedda aḍḍaṅki ayindi.} — The tree that fell on the road became a big obstacle.
  \item \emph{Recall:} say the Telugu for to take a turn, then the Telugu for a shortcut, then say \emph{aḍḍaṅki} again
\end{itemize}

\begin{tcolorbox}[breakable,colback=teal!4,colframe=teal!35!black,title={Before you move on}]
What does \te{అడ్డంకి} mean? (\textquotedblleft{}An obstacle; a hindrance\textquotedblright{}.) Say it once more.
\end{tcolorbox}
\section[\texorpdfstring{\te{మళ్ళింపు} (maḷḷimpu)}{మళ్ళింపు (maḷḷimpu)}]{\te{మళ్ళింపు} (maḷḷimpu) — a diversion; a detour}

\label{lesson:TE-C449-mallimpu}



\begin{quote}
Before the new one: say the Telugu for a shortcut, then the Telugu for an obstacle.
\end{quote}

\subsection*{You'll want to know: \te{మళ్ళింపు}}
\textbf{\te{మళ్ళింపు}} — \emph{maḷḷimpu} — \textquotedblleft{}a diversion; a detour\textquotedblright{}.

\begin{quote}
\te{వంతెన} \te{పనుల} \te{వల్ల} \te{రోడ్డు} \te{మీద} \te{మళ్ళింపు} \te{ఉంది}.
\end{quote}

\begin{quote}
\emph{vantena panula valla rōḍḍu mīda maḷḷimpu undi.}
\end{quote}

\begin{quote}
There is a diversion on the road because of work on the bridge.
\end{quote}

The noun of \te{మళ్ళించు}, to turn something aside. \te{మళ్ళీ}, again, comes from the same root: to turn back.

\begin{cousinweb}
\textbf{\te{మళ్ళింపు}} (\emph{maḷḷimpu}) is native Dravidian, the noun of \te{మళ్ళించు}, \textquotedblleft{}to turn aside\textquotedblright{}, from \te{మళ్ళు}, \textquotedblleft{}to turn back, to return\textquotedblright{} (DEDR 4761), the root of \te{మళ్ళీ}, \textquotedblleft{}again\textquotedblright{}.
\end{cousinweb}

\begin{grammarlens}[title={What you've built}]
One more word for getting around.
\end{grammarlens}

\subsection*{Your turn}
\begin{itemize}
\raggedright
  \item \emph{Say it:} \emph{maḷḷimpu}
  \item \emph{Say it:} \emph{maḷḷimpu}, once more
  \item \emph{Say it:} \emph{vantena panula valla rōḍḍu mīda maḷḷimpu undi.} — There is a diversion on the road because of work on the bridge.
  \item \emph{Recall:} say the Telugu for a shortcut, then the Telugu for an obstacle, then say \emph{maḷḷimpu} again
\end{itemize}

\begin{tcolorbox}[breakable,colback=teal!4,colframe=teal!35!black,title={Before you move on}]
What does \te{మళ్ళింపు} mean? (\textquotedblleft{}A diversion; a detour\textquotedblright{}.) Say it once more.
\end{tcolorbox}
\section[\texorpdfstring{\te{నిలిచిపోవు} (nilicipōvu)}{నిలిచిపోవు (nilicipōvu)}]{\te{నిలిచిపోవు} (nilicipōvu) — to stop dead; to come to a halt}

\label{lesson:TE-C449-nilicipovu}



\begin{quote}
Before the new one: say the Telugu for an obstacle, then the Telugu for a diversion.
\end{quote}

\subsection*{You'll want to know: \te{నిలిచిపోవు}}
\textbf{\te{నిలిచిపోవు}} — \emph{nilicipōvu} — \textquotedblleft{}to stop dead; to come to a halt\textquotedblright{}.

\begin{quote}
\te{వరద} \te{వల్ల} \te{రైళ్ళు} \te{నిలిచిపోయాయి}.
\end{quote}

\begin{quote}
\emph{varada valla raiḷḷu nilicipōyāyi.}
\end{quote}

\begin{quote}
Because of the flood the trains stopped running.
\end{quote}

Past (I …): \te{నిలిచిపోయాను}, \emph{nilicipōyānu}.

Said of traffic, a service or work that stops and stays stopped.

\begin{cousinweb}
\textbf{\te{నిలిచిపోవు}} (\emph{nilicipōvu}) is native Dravidian: \te{నిలుచు}, \textquotedblleft{}to stand, to stop, to stay\textquotedblright{} (DEDR 3675), with \te{పోవు}, \textquotedblleft{}to go\textquotedblright{} (DEDR 4572), which marks a change that has set in.
\end{cousinweb}

\begin{grammarlens}[title={What you've built}]
One more word for getting around.
\end{grammarlens}

\subsection*{Your turn}
\begin{itemize}
\raggedright
  \item \emph{Say it:} \emph{nilicipōvu}
  \item \emph{Say it:} \emph{nilicipōvu}, once more
  \item \emph{Say it:} \emph{varada valla raiḷḷu nilicipōyāyi.} — Because of the flood the trains stopped running.
  \item \emph{Recall:} say the Telugu for an obstacle, then the Telugu for a diversion, then say \emph{nilicipōvu} again
\end{itemize}

\begin{tcolorbox}[breakable,colback=teal!4,colframe=teal!35!black,title={Before you move on}]
What does \te{నిలిచిపోవు} mean? (\textquotedblleft{}To stop dead; to come to a halt\textquotedblright{}.) Say it once more.
\end{tcolorbox}
\section[\texorpdfstring{\te{స్తంభించు} (stambhiñcu)}{స్తంభించు (stambhiñcu)}]{\te{స్తంభించు} (stambhiñcu) — to seize up; to come to a standstill}

\label{lesson:TE-C449-stambhincu}



\begin{quote}
Before the new one: say the Telugu for a diversion, then the Telugu for to stop dead.
\end{quote}

\subsection*{You'll want to know: \te{స్తంభించు}}
\textbf{\te{స్తంభించు}} — \emph{stambhiñcu} — \textquotedblleft{}to seize up; to come to a standstill\textquotedblright{}.

\begin{quote}
\te{భారీ} \te{వర్షంతో} \te{నగరంలో} \te{ట్రాఫిక్} \te{స్తంభించింది}.
\end{quote}

\begin{quote}
\emph{bhārī varṣantō nagaraṁlō ṭrāphik stambhiñcindi.}
\end{quote}

\begin{quote}
With the heavy rain the traffic in the city came to a standstill.
\end{quote}

Past (it …): \te{స్తంభించింది}, \emph{stambhiñcindi}.

\te{స్తంభం} is a pillar: a road that stands still like a pillar. The news uses it of traffic and of normal life. Its subject is a thing, so the past shown is the third person: \te{స్తంభించింది}, it came to a standstill.

\begin{cousinweb}
\textbf{\te{స్తంభించు}} (\emph{stambhiñcu}) is built from Sanskrit \emph{stambh}, \textquotedblleft{}to become stiff, to be fixed\textquotedblright{}, whose noun \emph{stambha} is \textquotedblleft{}a pillar, a post\textquotedblright{} (CDIAL 13682), with the Telugu verb ending -\te{ఇంచు}.
\end{cousinweb}

\begin{grammarlens}[title={What you've built}]
One more word for getting around.
\end{grammarlens}

\subsection*{Your turn}
\begin{itemize}
\raggedright
  \item \emph{Say it:} \emph{stambhiñcu}
  \item \emph{Say it:} \emph{stambhiñcu}, once more
  \item \emph{Say it:} \emph{bhārī varṣantō nagaraṁlō ṭrāphik stambhiñcindi.} — With the heavy rain the traffic in the city came to a standstill.
  \item \emph{Recall:} say the Telugu for a diversion, then the Telugu for to stop dead, then say \emph{stambhiñcu} again
\end{itemize}

\begin{tcolorbox}[breakable,colback=teal!4,colframe=teal!35!black,title={Before you move on}]
What does \te{స్తంభించు} mean? (\textquotedblleft{}To seize up; to come to a standstill\textquotedblright{}.) Say it once more.
\end{tcolorbox}
\section[\texorpdfstring{\te{మూసుకుపోవు} (mūsukupōvu)}{మూసుకుపోవు (mūsukupōvu)}]{\te{మూసుకుపోవు} (mūsukupōvu) — to get blocked; to close up}

\label{lesson:TE-C449-musukupovu}



\begin{quote}
Before the new one: say the Telugu for to stop dead, then the Telugu for to seize up.
\end{quote}

\subsection*{You'll want to know: \te{మూసుకుపోవు}}
\textbf{\te{మూసుకుపోవు}} — \emph{mūsukupōvu} — \textquotedblleft{}to get blocked; to close up\textquotedblright{}.

\begin{quote}
\te{చెత్త} \te{వల్ల} \te{కాలువ} \te{మూసుకుపోయింది}.
\end{quote}

\begin{quote}
\emph{cetta valla kāluva mūsukupōyindi.}
\end{quote}

\begin{quote}
The canal got blocked with rubbish.
\end{quote}

Past (it …): \te{మూసుకుపోయింది}, \emph{mūsukupōyindi}.

\te{మూయు}, to close, + -\te{కో} + \te{పోవు}. A blocked nose is \te{ముక్కు} \te{మూసుకుపోయింది}. Its subject is a thing, so the past shown is the third person: \te{మూసుకుపోయింది}, it got blocked.

\begin{cousinweb}
\textbf{\te{మూసుకుపోవు}} (\emph{mūsukupōvu}) is native Dravidian: \te{మూయు}, \textquotedblleft{}to cover, to shut\textquotedblright{} (DEDR 4915), with the self-directed -\te{కో} and \te{పోవు}, \textquotedblleft{}to go\textquotedblright{} (DEDR 4572).
\end{cousinweb}

\begin{grammarlens}[title={What you've built}]
One more word for getting around.
\end{grammarlens}

\subsection*{Your turn}
\begin{itemize}
\raggedright
  \item \emph{Say it:} \emph{mūsukupōvu}
  \item \emph{Say it:} \emph{mūsukupōvu}, once more
  \item \emph{Say it:} \emph{cetta valla kāluva mūsukupōyindi.} — The canal got blocked with rubbish.
  \item \emph{Recall:} say the Telugu for to stop dead, then the Telugu for to seize up, then say \emph{mūsukupōvu} again
\end{itemize}

\begin{tcolorbox}[breakable,colback=teal!4,colframe=teal!35!black,title={Before you move on}]
What does \te{మూసుకుపోవు} mean? (\textquotedblleft{}To get blocked; to close up\textquotedblright{}.) Say it once more.
\end{tcolorbox}
